\documentclass[11pt]{article}
\usepackage[margin=0.75in]{geometry}
\usepackage{amsmath,amssymb}
\usepackage{graphicx}
\usepackage{kotex}
\usepackage{enumitem}
\usepackage{tikz}
\usetikzlibrary{arrows.meta,positioning,calc}

\setlength{\parindent}{0pt}
\setlength{\parskip}{0.2\baselineskip}
\renewcommand{\arraystretch}{1.12}
\setlist[enumerate]{topsep=1pt,itemsep=2pt,leftmargin=18pt}
\setlist[itemize]{topsep=1pt,itemsep=1pt,leftmargin=16pt}

\title{\vspace{-1.2em}3주차: 오차역전파 학습지}
\author{ML기초 시그 by 윤승현}
\date{}

\begin{document}
\maketitle
\vspace{-2.3em}

\section*{1. 계산 그래프로 보는 AND}

\[
z=x+y,\qquad s=z-\theta\;(\theta=1.5),\qquad \texttt{pred}=\mathbf{1}[s\ge 0]
\]

\begin{center}
\resizebox{0.88\textwidth}{!}{%
\begin{tikzpicture}[
>=Latex,
node distance=11mm and 13mm,
box/.style={draw, rounded corners=8pt, minimum height=9mm, minimum width=21mm, align=center, thick},
op/.style={box, fill=blue!18},
param/.style={box, fill=violet!18},
judge/.style={box, fill=green!18},
io/.style={box, fill=orange!18, minimum width=15mm}
]
\node[io] (x) {$x$};
\node[io, below=14mm of x] (y) {$y$};
\node[op, right=16mm of $(x)!0.5!(y)$, minimum width=28mm, minimum height=21mm] (sum) {\Large$+$};
\node[param, right=16mm of sum, minimum width=24mm] (sub) {$-\theta$};
\node[judge, right=16mm of sub, minimum width=26mm] (geq) {$0$ 이상};
\node[io, right=12mm of geq, fill=cyan!20] (pred) {\texttt{pred}};
\node[draw, rounded corners=5pt, fill=white, inner sep=2pt, below=1mm of sub] (theta) {$\theta=1.5$};

\draw[->, thick] (x.east) -- ++(5mm,0) |- (sum.west);
\draw[->, thick] (y.east) -- ++(5mm,0) |- (sum.west);
\draw[->, thick] (sum.east) -- (sub.west);
\draw[->, thick] (sub.east) -- (geq.west);
\draw[->, thick] (geq.east) -- (pred.west);
\end{tikzpicture}%
}
\end{center}

\begin{enumerate}
\item 가능한 모든 입력쌍 \((x,y)\in\{0,1\}^2\)에 대해 표를 완성하시오.

\begin{center}
\begin{tabular}{c c | c c c}
\hline
\(x\) & \(y\) & \(z=x+y\) & \(s=z-1.5\) & \texttt{pred} \\
\hline
0 & 0 & & & \\
0 & 1 & & & \\
1 & 0 & & & \\
1 & 1 & & & \\
\hline
\end{tabular}
\end{center}

\item \(Loss=(pred-true)^2\)라 하자. \((x,y)=(1,1)\), \(true=1\)일 때
\[
\frac{\partial Loss}{\partial \theta}
\]
를 구하시오.

\item
\[
\theta=\theta_{\mathrm{old}}-0.1\frac{\partial Loss}{\partial \theta}
\]
로 업데이트했을 때 Loss를 구하시오.

\item 위와 같이 업데이트했을 때 Loss가 줄어드는 이유를 간단히 설명하시오.
\end{enumerate}

\newpage

\section*{2. 손으로 훈련시켜봐요(XOR)}

\[
a=W_1x+b_1,\qquad
h=\mathrm{ReLU}(a),\qquad
\hat y=W_2h+b_2,\qquad
Loss=\frac12(\hat y-y)^2
\]

\[
x=
\begin{bmatrix}
1\\1
\end{bmatrix},\quad
W_1=
\begin{bmatrix}
1&0.5\\
0&1
\end{bmatrix},\quad
b_1=
\begin{bmatrix}
0\\0
\end{bmatrix}
\]

\[
W_2=
\begin{bmatrix}
1&1
\end{bmatrix},\qquad
b_2=0,\qquad
y=0
\]

\[
\mathrm{ReLU}(t)=\max(0,t)
\]

\begin{enumerate}
\item \textbf{순전파}
\[
a=\underline{\hspace{2.8cm}},\qquad
h=\underline{\hspace{2.8cm}},\qquad
\hat y=\underline{\hspace{1.8cm}},\qquad
Loss=\underline{\hspace{1.8cm}}
\]

\item \textbf{출력층 업데이트}

\[
\frac{\partial Loss}{\partial \hat y}=\hat y-y,\qquad
\frac{\partial \hat y}{\partial W_2}=h,\qquad
\frac{\partial \hat y}{\partial b_2}=1
\]

\[
\frac{\partial Loss}{\partial W_2}=\underline{\hspace{3cm}},\qquad
\frac{\partial Loss}{\partial b_2}=\underline{\hspace{2cm}}
\]

\[
W_2^{new}=W_2-0.1\frac{\partial Loss}{\partial W_2}
=\underline{\hspace{3.5cm}}
\]

\[
b_2^{new}=b_2-0.1\frac{\partial Loss}{\partial b_2}
=\underline{\hspace{2.5cm}}
\]

\item \textbf{ReLU 지나기}

\[
\mathrm{ReLU}'(t)=
\begin{cases}
1 & (t>0)\\
0 & (t<0)
\end{cases}
\qquad
\mathrm{ReLU}'(a)=\underline{\hspace{2.5cm}}
\]

\[
\frac{\partial Loss}{\partial a}
=\underline{\hspace{3cm}}
\]

\item \textbf{첫 층 업데이트}

\[
\frac{\partial Loss}{\partial W_1}
=\underline{\hspace{3.7cm}},\qquad
\frac{\partial Loss}{\partial b_1}
=\underline{\hspace{2.5cm}}
\]

\[
W_1^{new}=W_1-0.1\frac{\partial Loss}{\partial W_1}
=\underline{\hspace{3.5cm}}
\]

\[
b_1^{new}=b_1-0.1\frac{\partial Loss}{\partial b_1}
=\underline{\hspace{2.5cm}}
\]

\item \textbf{확인}

업데이트 후 다시 순전파하여
\[
\hat y_{\mathrm{new}}=\underline{\hspace{1.8cm}},\qquad
Loss_{\mathrm{new}}=\underline{\hspace{1.8cm}}
\]
를 구하시오.
\end{enumerate}

\end{document}